\documentclass[luatex,fontsize=8pt,paper=b5,tate]{jlreq}
\usepackage{luwa-ul,KKsymbols}
\usepackage{gckanbun,lltjext}

\begin{document}\fboxsep=0pt\fboxrule=.1pt%

\LARGE
\begin{GCKEnv}{2\zw}
  今夫\送り{レ}江戸\振り{者}{は}、世之所\送り{ノ}\返り{レ}称\送り{スル}名都\振り{大}{だい}\振り{邑}{いふ}、冠蓋之所\返り{レ}集\送り{マル}\LineNumbering*{\kakko{2}}\dashKKAuto{舟車之\振り{所}{ところ}\送り{ニシテ}\返り{レ}\振り{湊}{あつ}\送り[intrusion=post]{マル}}、実\送り{ニ}\振り{為}{た}\送り{ル}\返り{二}天下之大都会\返り{一}也。
  而\送り{レドモ}\LineNumbering{C}\underLineKKAuto{其地之為名、訪之於古、未之聞}。
  豈\送り{ニ}非\送り{ズ}\返り{三}古今相\送り{ヒ}去\送り{ルコト}\振り{日}{ひび}\送り{ニ}遠\送り{ク}、事之相\送り{ヒ}変\送り{ズルコト}愈多\送り{ク}、求\送り{ムルモ}\返り{二}其\送り{ノ}所\送り{ヲ}\返り{\IchiRe}欲\送り{スル}\返り{レ}聞\送り{カント}而不\送り{ルコト}可\送り{カラ}\返り{レ}得、亦\送り{タ}\振り{猶}{な}[ごと]\送り{ホ}[キヲ]\返り{二}今之於\送り{ケルガ}\返り{\JyouRe}古\送り{ニ}也。
\end{GCKEnv}

\bigskip

\begin{GCKEnv*}{2\zw}
  \fbox{\GCKanshiBox{10\zw}{%
    \振り[intrusion=both]{春}{しゅん}\hfill%
    \振り[intrusion=both]{眠}{みん}\hfill%
    不\返り[intrusion=both]{レ}\hfill%
    \振り{覚}{おぼ}\送り{エ}\返り[intrusion=both]{レ}\hfill%
    \振り{暁}{あかつき}\送り[intrusion=post]{ヲ}%
  }}

  \GCKanshiBox{10\zw}{%
    \振り[intrusion=both]{処}{しょ}\hfill%
    \振り[intrusion=both]{処}{しょ}\hfill%
    聞\送り{ク}\返り[intrusion=both]{二}\hfill%
    \underLineKKAuto{\振り[intrusion=both]{啼}{てい}\hfill%
    鳥\送り{ヲ}\返り[intrusion=post]{一}}%
  }

  \GCKanshiBox{10\zw}{%
    \振り[intrusion=both]{夜}{や}\hfill%
    \振り[intrusion=both]{来}{らい}\hfill%
    風\hfill%
    雨\送り[intrusion=both]{ノ}\hfill%
    声%
  }

  \GCKanshiBox{10\zw}{%
    花\hfill%
    \振り{落}{お}\送り[intrusion=both]{ツルコト}\hfill%
    知\送り[intrusion=both]{ル}\hfill%
    \振り[intrusion=both]{多}{た}\hfill%
    \振り[intrusion=both]{少}{しょう}%
  }
\end{GCKEnv*}

\end{document}